\section*{Capítulo \ref{ch:g10:muscles-and-joints} --- Músculos y articulaciones}

\begin{solution}{exo:g10:muscles-and-joints:1}
\omterm{def:g10:muscles-and-joints:joint}{Cartílago} en los extremos de los huesos (deslizamiento suave, reparto
de la carga); cápsula (sella la \omterm{def:g10:muscles-and-joints:joint}{articulación}) con su líquido sinovial
(lubricación); \omterm{def:g10:muscles-and-joints:joint}{ligamentos} (mantienen unidos los huesos y limitan los
movimientos); en la rodilla, meniscos (reparten la carga).
\end{solution}

\begin{solution}{exo:g10:muscles-and-joints:2}
Un \omterm{def:g10:muscles-and-joints:muscle}{tendón} une un músculo a un hueso y transmite su tirón; un \omterm{def:g10:muscles-and-joints:joint}{ligamento}
une hueso con hueso cruzando una \omterm{def:g10:muscles-and-joints:joint}{articulación} y limita su movimiento.
Ninguno de los dos se contrae.
\end{solution}

\begin{solution}{exo:g10:muscles-and-joints:3}
Un músculo solo puede tirar acortándose; no puede empujar. Para mover
una \omterm{def:g10:muscles-and-joints:joint}{articulación} en los dos sentidos hacen falta un flexor y un
extensor, cada uno de los cuales deshace lo que hizo el otro.
\end{solution}

\begin{solution}{exo:g10:muscles-and-joints:4}
Esguince: un \omterm{def:g10:muscles-and-joints:joint}{ligamento} estirado o roto. Rotura muscular: \omterm{def:g10:muscles-and-joints:muscle}{fibras
musculares} desgarradas. Luxación: el extremo de un hueso forzado fuera
de su \omterm{def:g10:muscles-and-joints:joint}{articulación}, con la cápsula y los \omterm{def:g10:muscles-and-joints:joint}{ligamentos} estirados o rotos.
\end{solution}

\begin{solution}{exo:g10:muscles-and-joints:5}
Fibras lentas: lentas, poco potentes, incansables, ricas en
\omterm{def:g10:cells-common-unit:organelle}{mitocondrias} y capilares, rojas, funcionan con la respiración. Fibras
rápidas: rápidas, potentes, se cansan en un minuto, pocas \omterm{def:g10:cells-common-unit:organelle}{mitocondrias},
pálidas, funcionan con la reserva fosforada y la \omterm{def:g10:cell-metabolism:fermentation}{fermentación}.
\end{solution}

\begin{solution}{exo:g10:muscles-and-joints:6}
Carga $2 \times 9.8 = \qty{19.6}{N}$; fuerza del bíceps $19.6 \times 32/4
\approx \qty{157}{N}$.
\end{solution}

\begin{solution}{exo:g10:muscles-and-joints:7}
El empujón puede venir en cualquier dirección, y cada músculo solo
puede resistir una. Contraerse los dos rigidiza la \omterm{def:g10:muscles-and-joints:joint}{articulación} en
ambos sentidos a la vez; venga por donde venga el empujón, uno de ellos
ya está tirando en contra.
\end{solution}

\begin{solution}{exo:g10:muscles-and-joints:8}
La fibra se acorta \qty{0.6}{cm}; el \omterm{def:g10:muscles-and-joints:muscle}{tendón} mueve el hueso
\qty{0.6}{cm} a \qty{5}{cm} del codo, así que la mano se desplaza
$0.6 \times 35/5 = \qty{4.2}{cm}$.
\end{solution}

\begin{solution}{exo:g10:muscles-and-joints:9}
La movilidad exige una cavidad poco profunda y \omterm{def:g10:muscles-and-joints:joint}{ligamentos} laxos; esa
misma laxitud deja que la cabeza del húmero salga de la cavidad cuando
una fuerza la empuja más allá del recorrido que los \omterm{def:g10:muscles-and-joints:joint}{ligamentos} pueden
sujetar.
\end{solution}

\begin{solution}{exo:g10:muscles-and-joints:10}
La deshidratación y la pérdida de sales por el sudor, y la fatiga de
las fibras al final de la carrera. Remedio: estirar el músculo con
suavidad y beber agua con sal.
\end{solution}

\begin{solution}{exo:g10:muscles-and-joints:11}
El muslo del velocista es un 70\% de fibras rápidas, y la proporción
está en buena medida heredada; el entrenamiento agranda las fibras y
mejora su equipamiento, pero no convierte su número. La persona no
entrenada parte de mitad y mitad y tiene más fibras lentas que
desarrollar.
\end{solution}

\begin{solution}{exo:g10:muscles-and-joints:12}
$4 \times 70 \times 9.8 \approx \qty{2.7}{kN}$. Tensión $2744/100
\approx \qty{27}{N/mm^2}$: alrededor de una cuarta parte de la tensión
de rotura, un factor de seguridad de cuatro.
\end{solution}

\begin{solution}{exo:g10:muscles-and-joints:13}
El músculo está muy irrigado y contiene \omterm{def:g10:cells-common-unit:cell}{células} de reserva que
reconstruyen fibras: semanas. Los \omterm{def:g10:muscles-and-joints:joint}{ligamentos} tienen pocos vasos, así
que los nutrientes y las \omterm{def:g10:cells-common-unit:cell}{células} reparadoras llegan despacio: meses. El
\omterm{def:g10:muscles-and-joints:joint}{cartílago} no tiene ningún vaso sanguíneo y casi ninguna \omterm{def:g10:cells-common-unit:cell}{célula} que se
divida; se alimenta solo del líquido sinovial y apenas se repara.
\end{solution}

\begin{solution}{exo:g10:muscles-and-joints:14}
El culturismo engrosa las fibras (más miofibrillas en cada una) y los
estiramientos alargan el conjunto músculo--tendón y la tolerancia de la
cápsula. El número de fibras queda fijado al principio de la vida; el
músculo adulto crece o mengua cambiando el tamaño de las fibras que ya
tiene.
\end{solution}

\begin{solution}{exo:g10:muscles-and-joints:15}
En las escaleras, la \omterm{def:g10:muscles-and-joints:joint}{articulación} soporta varias veces el peso
corporal; un \omterm{def:g10:muscles-and-joints:joint}{cartílago} más fino transmite esa presión al hueso de
debajo, que tiene nervios y duele. Con menos \omterm{def:g10:muscles-and-joints:joint}{cartílago}, los meniscos
son lo último que reparte la carga. Unos músculos fuertes alrededor de
la rodilla absorben parte del impacto y estabilizan la \omterm{def:g10:muscles-and-joints:joint}{articulación}, y
así reducen la presión máxima sobre el \omterm{def:g10:muscles-and-joints:joint}{cartílago} que queda.
\end{solution}

\begin{solution}{pb:g10:muscles-and-joints:1}
\textbf{1.} $70 \times 9.8 = \qty{686}{N}$.

\textbf{2.} $686 \times 10/5 \approx \qty{1.4}{kN}$.

\textbf{3.} $686 \times 20/5 \approx \qty{2.7}{kN}$.

\textbf{4.} Dos y cuatro veces el peso corporal: en cada escalón, el
cuádriceps tira con el peso de cuatro personas, y por eso arden los
muslos en una subida larga.

\textbf{5.} Los isquiotibiales, en la parte posterior del muslo; se
relajan y se estiran mientras el cuádriceps estira la rodilla, y se
contraen ligeramente para estabilizar la \omterm{def:g10:muscles-and-joints:joint}{articulación}.

\textbf{6.} Fuerza del suelo $3 \times 686 \approx \qty{2.06}{kN}$;
fuerza del \omterm{def:g10:muscles-and-joints:muscle}{tendón} $2058 \times 20/5 \approx \qty{8.2}{kN}$, doce veces
el peso corporal.

\textbf{7.} $8232/120 \approx \qty{69}{N/mm^2}$: unas dos terceras
partes de la tensión de rotura. Cada aterrizaje consume casi todo el
margen del \omterm{def:g10:muscles-and-joints:muscle}{tendón}.

\textbf{8.} Carga $8.2 + 2.1 \approx \qty{10.3}{kN}$. Con los meniscos:
$10300/600 \approx \qty{17}{N/mm^2}$; sin ellos: $10300/100 \approx
\qty{103}{N/mm^2}$, seis veces más.

\textbf{9.} Seis veces la presión sobre el mismo \omterm{def:g10:muscles-and-joints:joint}{cartílago} en cada
aterrizaje: la superficie se machaca más deprisa de lo que puede
mantenerse, y se desgasta.

\textbf{10.} Una película de líquido entre los \omterm{def:g10:muscles-and-joints:joint}{cartílagos}, para que
deslicen durante la flexión sin rozamiento seco, y una ligera
amortiguación del impacto.

\textbf{11.} Un giro de la tibia bajo el fémur, con una flexión
lateral: rotaciones para las que la rodilla no está hecha y a las que
los \omterm{def:g10:muscles-and-joints:joint}{ligamentos} están ahí para resistirse.

\textbf{12.} El cuádriceps solo tira a lo largo de la pierna y estira
la rodilla; no tiene ninguna acción contra un giro. Y la lesión ocurre
en unas centésimas de segundo, más deprisa que cualquier respuesta
muscular.

\textbf{13.} Sangre de los vasos del \omterm{def:g10:muscles-and-joints:joint}{ligamento} roto y de la cápsula,
más líquido sinovial de más: la cápsula quedó rota o inflamada y la
\omterm{def:g10:muscles-and-joints:joint}{articulación} se llenó.

\textbf{14.} El \omterm{def:g10:muscles-and-joints:joint}{ligamento}, mal irrigado, necesita meses y a menudo
cirugía; el menisco desgarrado, \omterm{def:g10:muscles-and-joints:joint}{cartílago} casi sin vasos, no suele
curar y la parte rota se recorta o se sutura.

\textbf{15.} Los isquiotibiales son los antagonistas del cuádriceps y
tiran de la tibia hacia atrás; una pareja fuerte que se contraiga a la
vez rigidiza la rodilla frente al deslizamiento hacia delante que el
\omterm{def:g10:muscles-and-joints:joint}{ligamento} roto ya no impide, y protege la \omterm{def:g10:muscles-and-joints:joint}{articulación} antes y después
de la cirugía.

\textbf{16.} Unas 4000 veces al día a \qty{1.4}{kN} o más.

\textbf{17.} Unas $\num{1.5e6}$ cargas al año. La tendinitis es un daño
que se acumula más deprisa que la lenta reparación del tejido
tendinoso; solo el reposo deja que la reparación alcance al daño, ya
que la fuerza no reduce la carga de cada paso.

\textbf{18.} Aterrizajes: $40 \times 8.2 \approx \qty{330}{kN}$ en
total; pasos: $4000 \times 1.4 \approx \qty{5600}{kN}$ --- los pasos
soportan diecisiete veces más carga total, pero los aterrizajes
soportan los picos que amenazan con romper.

\textbf{19.} Duplicar la fuerza del suelo duplica las fuerzas del
\omterm{def:g10:muscles-and-joints:muscle}{tendón} y de la \omterm{def:g10:muscles-and-joints:joint}{articulación} en cada zancada; la bicicleta y la natación
trabajan los músculos cargando la rodilla muy por debajo del peso
corporal.

\textbf{20.} El \omterm{def:g10:muscles-and-joints:muscle}{tendón} del cuádriceps soporta cuatro veces el peso
corporal en un escalón y unas doce en un aterrizaje; los \omterm{def:g10:muscles-and-joints:joint}{ligamentos}, a
los que esas cargas nunca ponen a prueba, fueron el tejido que encontró
el giro.
\end{solution}
